% ======================================================================
\section{Solver verification, cost, and scaling}
\label{sec:res_scaling}
% ======================================================================

The production distributed solver is validated against the dense reference and
benchmarked on the two axes that matter for anharmonic transport: the
block-tridiagonal recursion along the device and the energy-point parallelism
across MPI ranks. The two solvers follow the structure of the self-energy
(\cref{ssec:phph_distribution,ssec:phph_qdist}): Gamma-only conductors ($k{=}1$,
the nanowires and tubes) scale on $\texttt{block\_comm\_size}\times\texttt{stack}$,
while the transversely-periodic systems ($k{>}1$, films and perovskites) scale on
$\texttt{q\_comm\_size}\times\texttt{stack}$ with $\texttt{block\_comm\_size}{=}1$;
for the large-block Gamma-only regime the inner-ring contraction additionally
threads over the frequency axis (\cref{ssec:phph_ring}).

\paragraph{Machine-precision agreement with the dense reference.} The production
block-tridiagonal recursive Green's-function solver
(\texttt{qttools.greens\_function\_solver.RGF}) reproduces the dense reference and a
direct \texttt{numpy} inverse: the generic selected-inversion suite passes all
$144$ parametrised cases, and an end-to-end check on a genuine phonon Dyson matrix
$A(\omega^2)=\omega^2 I-H_D-\Sigma_\mathrm{leads}$ with Sancho--Rubio
open-boundary self-energies agrees to $1.4\times10^{-16}$ on the selected retarded
Green's function across $\omega\in[1,16]$\,THz. The three-phonon bubble and
scattering-self-energy kernels pass all $28$ tests; these also flag that a
nearest-neighbour FC3 truncation discards $34$--$48\%$ of the Frobenius norm for the
test fixtures, so the number of primitive cells per transport cell must be chosen
with care. The block-distributed recursion
(\texttt{rgf\_dist.RGFDist}) reproduces the serial reference under
\texttt{mpirun -np 3} ($\texttt{block\_comm\_size}{=}3$).

\paragraph{Solver cost.} The dense reference inverts the full $N_D\times N_D$ device
matrix, so its cost grows as $\mathcal{O}(n_\mathrm{slabs}^3)$, whereas RGF exploits
the block-tridiagonal structure and is linear in the number of cells
(\cref{fig:res_solver_scaling}, left): on phonon-Dyson matrices with a fixed
$192$-dof block, from $2$ to $32$ blocks the RGF wall time rises linearly while the
dense time follows the cubic guide, a $58\times$ gap at $32$ blocks, with the two
agreeing to $\sim10^{-13}$ on the block-tridiagonal entries. The complementary axis
-- distributing the energy points across MPI ranks -- is shown on the right: a fixed
$1536\times1536$ problem with $128$ energy points falls from $1.55$\,s on one rank to
$0.86$, $0.55$ and $0.28$\,s on $2$, $4$ and $8$ ranks, a $5.3$--$5.5\times$ speed-up
at eight ranks ($68\%$ parallel efficiency) that improves with problem size.

\begin{figure}[htb]
  \centering
  \includegraphics[width=0.48\linewidth]{transport_sweeps/solver_scaling.pdf}
  \caption{Solver scaling. Left: selected-inversion wall time vs.\ number of
  transport-cell blocks -- RGF is linear, the dense reference cubic ($58\times$
  apart at $32$ blocks), agreeing to $\sim10^{-13}$. Right: distributed
  energy-parallel speed-up ($5.3\times$ at eight ranks).}
  \label{fig:res_solver_scaling}
\end{figure}

This separation is decisive for the anharmonic problem. The dense self-energy bubble
parallelises across the $(I,J)$ slab-pair blocks -- their count grows
$1\to8\to15$ for $n_\mathrm{slabs}=1,2,3$ -- but each block's bubble is itself a
serial BLAS contraction, so the dense per-point wall time climbs steeply: a single
$(3,3)$-CNT SCBA point costs $66$\,s at $n_\mathrm{slabs}{=}1$ and $1193$\,s
($\sim18\times$) at $n_\mathrm{slabs}{=}2$, and the converged three-cell point
($108$ DOF, $15$ blocks) costs $6.8$\,h. Multi-cell anharmonic transport is therefore
the regime of the RGF and energy-parallel production path, with the dense solver kept
as the verification oracle.

\todo[inline]{Add the production phph throughput study (the restored
\texttt{phph\_scaling}/\texttt{phph\_q\_scaling}/\texttt{phph\_memory}/%
\texttt{prod\_phph\_scaling} figures in \texttt{fig/transport\_sweeps/}): strong/weak
scaling on $\texttt{block}\times\texttt{stack}$ and $\texttt{q}\times\texttt{stack}$,
the ring-thread throughput curve, and host/GPU memory vs.\ system size. See
\texttt{document/MISSING.md} (A2).}
